% GENERATED FILE. Edit canonical lessons, then run npm run generate:books.
% canonical-source-hash: fnv1a64:c5c35541c9fc3c4f
% canonical-lessons: BN-C410-math, BN-C410-park, BN-C410-uthon, BN-C410-chotvor, BN-C410-paychari-kora

\chapter{Fields and Parks}
\label{ch:bn-fields-and-parks}
\hlchaptermodality{410}

\begin{chapteropening}
\textbf{By the end of this chapter:} \emph{I can talk about the playing field, the park, the courtyard and the square.}

Complete the last lesson of chapter 410: I can talk about the playing field, the park, the courtyard and the square.
\end{chapteropening}

\section[\texorpdfstring{\bn{মাঠ} (māṭh)}{মাঠ (māṭh)}]{\bn{মাঠ} (māṭh) — a field; an open ground}

\label{lesson:BN-C410-math}



\begin{quote}
Before the new one: say the Bengali for a resident, then the Bengali for to hang out (chatting).
\end{quote}

\subsection*{You'll want to know: \bn{মাঠ}}
\textbf{\bn{মাঠ}} — \emph{māṭh} — \textquotedblleft{}a field; an open ground\textquotedblright{}.

\begin{quote}
\bn{বিকেলে} \bn{ছেলেরা} \bn{মাঠে} \bn{ফুটবল} \bn{খেলে।}
\end{quote}

\begin{quote}
\emph{bikele chhelerā māṭhe phuṭbôl khele.}
\end{quote}

\begin{quote}
In the afternoons the boys play football on the field.
\end{quote}

Open ground for games, meetings and grazing; a field of crops is a \bn{খেত}.

\begin{cousinweb}
\textbf{\bn{মাঠ}} (\emph{māṭh}) is the open ground of every town and village; where the word comes from is not settled, and CDIAL gives no Bengali form for it.
\end{cousinweb}

\begin{grammarlens}[title={What you've built}]
One more word for getting around.
\end{grammarlens}

\subsection*{Your turn}
\begin{itemize}
\raggedright
  \item \emph{Say it:} \emph{māṭh}
  \item \emph{Say it:} \emph{māṭh}, once more
  \item \emph{Say it:} \emph{bikele chhelerā māṭhe phuṭbôl khele.} — In the afternoons the boys play football on the field.
  \item \emph{Recall:} say the Bengali for a resident, then the Bengali for to hang out (chatting), then say \emph{māṭh} again
\end{itemize}

\begin{tcolorbox}[breakable,colback=teal!4,colframe=teal!35!black,title={Before you move on}]
What does \bn{মাঠ} mean? (\textquotedblleft{}A field; an open ground\textquotedblright{}.) Say it once more.
\end{tcolorbox}
\section[\texorpdfstring{\bn{পার্ক} (pārk)}{পার্ক (pārk)}]{\bn{পার্ক} (pārk) — a park}

\label{lesson:BN-C410-park}



\begin{quote}
Before the new one: say the Bengali for to hang out (chatting), then the Bengali for an open ground.
\end{quote}

\subsection*{You'll want to know: \bn{পার্ক}}
\textbf{\bn{পার্ক}} — \emph{pārk} — \textquotedblleft{}a park\textquotedblright{}.

\begin{quote}
\bn{সকালে} \bn{বুড়ো} \bn{লোকেরা} \bn{পার্কে} \bn{হাঁটেন।}
\end{quote}

\begin{quote}
\emph{sôkāle buṛo lokerā pārke hā̃ṭen.}
\end{quote}

\begin{quote}
In the mornings old people walk in the park.
\end{quote}

\begin{cousinweb}
\textbf{\bn{পার্ক}} (\emph{pārk}) is borrowed from English \emph{park}.
\end{cousinweb}

\begin{grammarlens}[title={What you've built}]
One more word for getting around.
\end{grammarlens}

\subsection*{Your turn}
\begin{itemize}
\raggedright
  \item \emph{Say it:} \emph{pārk}
  \item \emph{Say it:} \emph{pārk}, once more
  \item \emph{Say it:} \emph{sôkāle buṛo lokerā pārke hā̃ṭen.} — In the mornings old people walk in the park.
  \item \emph{Recall:} say the Bengali for to hang out (chatting), then the Bengali for an open ground, then say \emph{pārk} again
\end{itemize}

\begin{tcolorbox}[breakable,colback=teal!4,colframe=teal!35!black,title={Before you move on}]
What does \bn{পার্ক} mean? (\textquotedblleft{}A park\textquotedblright{}.) Say it once more.
\end{tcolorbox}
\section[\texorpdfstring{\bn{উঠোন} (uṭhon)}{উঠোন (uṭhon)}]{\bn{উঠোন} (uṭhon) — a courtyard (of a house)}

\label{lesson:BN-C410-uthon}



\begin{quote}
Before the new one: say the Bengali for an open ground, then the Bengali for a park.
\end{quote}

\subsection*{You'll want to know: \bn{উঠোন}}
\textbf{\bn{উঠোন}} — \emph{uṭhon} — \textquotedblleft{}a courtyard (of a house)\textquotedblright{}.

\begin{quote}
\bn{দিদিমা} \bn{উঠোনে} \bn{বসে} \bn{চাল} \bn{বাছছিলেন।}
\end{quote}

\begin{quote}
\emph{didimā uṭhone bôse chāl bāchhchhilen.}
\end{quote}

\begin{quote}
Grandmother sat in the courtyard picking over the rice.
\end{quote}

The open yard in the middle of a village house, where work is done and children play.

\begin{cousinweb}
\textbf{\bn{উঠোন}} (\emph{uṭhon}) is inherited: CDIAL lists Bengali \emph{uṭhān}, \textquotedblleft{}a courtyard\textquotedblright{}, under Sanskrit \emph{upasth\'{\={a}}na}, \textquotedblleft{}approach, attendance\textquotedblright{} and later \textquotedblleft{}an abode\textquotedblright{} (CDIAL 2275); \bn{উঠোন} is the everyday spoken form.
\end{cousinweb}

\begin{grammarlens}[title={What you've built}]
One more word for getting around.
\end{grammarlens}

\subsection*{Your turn}
\begin{itemize}
\raggedright
  \item \emph{Say it:} \emph{uṭhon}
  \item \emph{Say it:} \emph{uṭhon}, once more
  \item \emph{Say it:} \emph{didimā uṭhone bôse chāl bāchhchhilen.} — Grandmother sat in the courtyard picking over the rice.
  \item \emph{Recall:} say the Bengali for an open ground, then the Bengali for a park, then say \emph{uṭhon} again
\end{itemize}

\begin{tcolorbox}[breakable,colback=teal!4,colframe=teal!35!black,title={Before you move on}]
What does \bn{উঠোন} mean? (\textquotedblleft{}A courtyard (of a house)\textquotedblright{}.) Say it once more.
\end{tcolorbox}
\section[\texorpdfstring{\bn{চত্বর} (chôtvôr)}{চত্বর (chôtvôr)}]{\bn{চত্বর} (chôtvôr) — a square; an open paved area}

\label{lesson:BN-C410-chotvor}



\begin{quote}
Before the new one: say the Bengali for a park, then the Bengali for a courtyard.
\end{quote}

\subsection*{You'll want to know: \bn{চত্বর}}
\textbf{\bn{চত্বর}} — \emph{chôtvôr} — \textquotedblleft{}a square; an open paved area\textquotedblright{}.

\begin{quote}
\bn{মন্দিরের} \bn{চত্বরে} \bn{শিশুরা} \bn{খেলছিল।}
\end{quote}

\begin{quote}
\emph{mondirer chôtvôre shishurā khelchhilô.}
\end{quote}

\begin{quote}
Children were playing in the temple square.
\end{quote}

\begin{cousinweb}
\textbf{\bn{চত্বর}} (\emph{chôtvôr}) is Sanskrit \emph{catvara}, \textquotedblleft{}a quadrangular place, a place where many ways meet\textquotedblright{} (Monier-Williams), taken over whole.
\end{cousinweb}

\begin{grammarlens}[title={What you've built}]
One more word for getting around.
\end{grammarlens}

\subsection*{Your turn}
\begin{itemize}
\raggedright
  \item \emph{Say it:} \emph{chôtvôr}
  \item \emph{Say it:} \emph{chôtvôr}, once more
  \item \emph{Say it:} \emph{mondirer chôtvôre shishurā khelchhilô.} — Children were playing in the temple square.
  \item \emph{Recall:} say the Bengali for a park, then the Bengali for a courtyard, then say \emph{chôtvôr} again
\end{itemize}

\begin{tcolorbox}[breakable,colback=teal!4,colframe=teal!35!black,title={Before you move on}]
What does \bn{চত্বর} mean? (\textquotedblleft{}A square; an open paved area\textquotedblright{}.) Say it once more.
\end{tcolorbox}
\section[\texorpdfstring{\bn{পায়চারি} \bn{করা} (pāychāri kôrā)}{পায়চারি করা (pāychāri kôrā)}]{\bn{পায়চারি} \bn{করা} (pāychāri kôrā) — to pace up and down; to stroll}

\label{lesson:BN-C410-paychari-kora}



\begin{quote}
Before the new one: say the Bengali for a courtyard, then the Bengali for a square (open place).
\end{quote}

\subsection*{You'll want to know: \bn{পায়চারি} \bn{করা}}
\textbf{\bn{পায়চারি} \bn{করা}} — \emph{pāychāri kôrā} — \textquotedblleft{}to pace up and down; to stroll\textquotedblright{}.

\begin{quote}
\bn{অরুণ} \bn{রাতে} \bn{ছাদে} \bn{পায়চারি} \bn{করল।}
\end{quote}

\begin{quote}
\emph{ôrun rāte chhāde pāychāri kôrlô.}
\end{quote}

\begin{quote}
At night Arun paced up and down on the roof.
\end{quote}

Past (he/she …): \bn{তিনি} \bn{বারান্দায়} \bn{পায়চারি} \bn{করলেন}, \emph{tini bārāndāy pāychāri kôrlen}.

\begin{cousinweb}
\textbf{\bn{পায়চারি} \bn{করা}} (\emph{pāychāri kôrā}) joins \bn{পায়চারি}, \textquotedblleft{}walking to and fro\textquotedblright{}, to \bn{করা}, to do; the origin of \bn{পায়চারি} is not settled.
\end{cousinweb}

\begin{grammarlens}[title={What you've built}]
One more word for getting around.
\end{grammarlens}

\subsection*{Your turn}
\begin{itemize}
\raggedright
  \item \emph{Say it:} \emph{pāychāri kôrā}
  \item \emph{Say it:} \emph{pāychāri kôrā}, once more
  \item \emph{Say it:} \emph{ôrun rāte chhāde pāychāri kôrlô.} — At night Arun paced up and down on the roof.
  \item \emph{Recall:} say the Bengali for a courtyard, then the Bengali for a square (open place), then say \emph{pāychāri kôrā} again
\end{itemize}

\begin{tcolorbox}[breakable,colback=teal!4,colframe=teal!35!black,title={Before you move on}]
What does \bn{পায়চারি} \bn{করা} mean? (\textquotedblleft{}To pace up and down; to stroll\textquotedblright{}.) Say it once more.
\end{tcolorbox}
